\section{Equalizatory, koequalizatory i kwadraty kartezjańskie}\label{sec:equalizatory}

\providecommand{\atlasref}[2]{#2}

W konstrukcji iloczynu porównujemy dwie strzałki o wspólnej dziedzinie. Czasem potrzebujemy z kolei \emph{wymusić równość} dwóch strzałek lub połączyć dwa obiekty nad wspólną bazą. Te potrzeby prowadzą do equalizatorów i kwadratów kartezjańskich.

\subsection{Equalizator i koequalizator}

\begin{definicja}[Equalizator]\label{def:equalizer}
Niech $f,g:A\to B$ będą dwoma morfizmami w kategorii $\mathcal C$. \emph{Equalizatorem} $f$ i $g$ nazywamy morfizm $e:E\to A$ taki, że $f\circ e=g\circ e$, oraz taki, że dla każdego $u:X\to A$ z $f\circ u=g\circ u$ istnieje dokładnie jeden $\bar u:X\to E$ spełniający $e\circ\bar u=u$.
\end{definicja}
Dla $fu=gu$ własność uniwersalną przedstawia diagram:
\begin{center}
\begin{tikzpicture}[>=stealth,every node/.style={font=\small}]
 \node (x) at (0,2) {$X$};
 \node (e) at (0,0) {$E$};
 \node (a) at (3.4,0) {$A$};
 \node (b) at (7,0) {$B$};
 \draw[->] (x) -- node[above right] {$u$} (a);
 \draw[->,dashed] (x) -- node[left] {$\bar u$} (e);
 \draw[{Hooks[right]}-Stealth] (e) -- node[below] {$e$} (a);
 \draw[->] (a) to[bend left=18] node[above] {$f$} (b);
 \draw[->] (a) to[bend right=18] node[below] {$g$} (b);
\end{tikzpicture}
\end{center}

\begin{przyklad}[Equalizator w zbiorach]
Dla funkcji $f,g:A\to B$ przyjmijmy $E=\{a\in A:f(a)=g(a)\}$, a za $e:E\hookrightarrow A$ weźmy włączenie. Jeśli $u:X\to A$ spełnia $fu=gu$, to $u(x)\in E$ dla każdego $x$, zatem $u$ ma dokładnie jedno rozłożenie przez $e$. Na przykład dla $f,g:\mathbb R\to\mathbb R$ danych przez $f(x)=x^2$ i $g(x)=x$ otrzymujemy $E=\{0,1\}$.
\end{przyklad}

\begin{przyklad}[Jądro jako equalizator]
W $\Vect_k$ jądro odwzorowania $f:V\to W$ jest equalizatorem pary $f,0:V\to W$, gdzie $0$ jest odwzorowaniem zerowym. Istotnie, $f|_{\Ker f}=0$, a każde $u:X\to V$ z $fu=0$ przyjmuje wartości w $\Ker f$. To pierwsza wskazówka, że jądro jest konstrukcją kategorialną, a nie wyłącznie zbiorem rozwiązań równania.
\end{przyklad}

\begin{definicja}[Koequalizator]\label{def:coeq}
Dla morfizmów $f,g:A\to B$ \emph{koequalizatorem} nazywamy morfizm $q:B\to Q$ taki, że $q\circ f=q\circ g$ oraz dla każdego $v:B\to X$ spełniającego $v\circ f=v\circ g$ istnieje dokładnie jeden $\bar v:Q\to X$ o własności $\bar v\circ q=v$.
\end{definicja}
Dla $vf=vg$ diagram dualny ma postać:
\begin{center}
\begin{tikzpicture}[>=stealth,every node/.style={font=\small}]
 \node (a) at (0,0) {$A$};
 \node (b) at (3.6,0) {$B$};
 \node (q) at (7,0) {$Q$};
 \node (x) at (7,2) {$X$};
 \draw[->] (a) to[bend left=18] node[above] {$f$} (b);
 \draw[->] (a) to[bend right=18] node[below] {$g$} (b);
 \draw[->>] (b) -- node[below] {$q$} (q);
 \draw[->] (b) -- node[above left] {$v$} (x);
 \draw[->,dashed] (q) -- node[right] {$\bar v$} (x);
\end{tikzpicture}
\end{center}

\begin{przyklad}[Iloraz zbioru]
W $\Set$ koequalizator można zbudować, biorąc najmniejszą relację równoważności $\sim$ na $B$, dla której $f(a)\sim g(a)$ dla każdego $a\in A$, oraz rzut $q:B\to B/{\sim}$. Relacja $\sim$ istnieje jako przecięcie wszystkich relacji równoważności zawierających wskazane pary. Dla funkcji $v:B\to X$ spełniającej $vf=vg$ relacja $b\equiv b'$ określona przez $v(b)=v(b')$ jest relacją równoważności zawierającą te pary. Zatem $\sim$ zawiera się w $\equiv$, a wzór $\bar v([b])=v(b)$ jest dobrze określony. Surjektywność $q$ wymusza jednoznaczność.

Jeśli $A=\{*\}$, $B=\{0,1,2\}$, $f(*)=0$ i $g(*)=1$, to $Q$ ma dwie klasy: $\{0,1\}$ i $\{2\}$. Koequalizator \emph{skleja} obrazy wskazane przez parę morfizmów.
\end{przyklad}

\begin{przyklad}[Kojądro]
W $\Vect_k$ \emph{kojądro} odwzorowania $f:V\to W$ jest koequalizatorem pary $f,0:V\to W$. Ma postać rzutu $q:W\to W/\im f$. Jeśli $v:W\to X$ jest liniowa i $vf=0$, to $\im f\subseteq\Ker v$, więc $v$ przechodzi jednoznacznie przez $q$.
\end{przyklad}

\begin{stwierdzenie}
Każdy equalizator jest monomorfizmem. Każdy koequalizator jest epimorfizmem.
\end{stwierdzenie}
\begin{proof}
Niech $e:E\to A$ będzie equalizatorem $f,g:A\to B$. Jeśli $eu=ev$ dla $u,v:X\to E$, to wspólny morfizm $w=eu=ev$ spełnia $fw=gw$, więc na mocy jedyności rozkładu przez $e$ mamy $u=v$. Jeżeli $q:B\to Q$ jest koequalizatorem, a $u,v:Q\to X$ spełniają $u\circ q=v\circ q$, to wspólny morfizm $w=u\circ q$ spełnia $w\circ f=w\circ g$. Jednoznaczność faktoryzacji przez $q$ daje $u=v$, więc $q$ jest epimorfizmem.
\end{proof}


\begin{przyklad}[Equalizator dwóch homomorfizmów]
W kategorii $\Ab$ rozważmy $f,g:\mathbb Z\to\mathbb Z/6\mathbb Z$ określone przez $f(n)=[2n]_6$ oraz $g(n)=[5n]_6$. Warunek $f(n)=g(n)$ jest równoważny $3n\equiv0\pmod 6$, a więc $n$ jest parzyste. Equalizatorem jest zanurzenie $2\mathbb Z\hookrightarrow\mathbb Z$. Jeśli homomorfizm $u:H\to\mathbb Z$ spełnia $fu=gu$, to jego obraz jest zawarty w $2\mathbb Z$; dokładnie stąd wynika istnienie i jedyność wymaganego rozkładu. To przykład, w którym definicja uniwersalna daje równanie arytmetyczne.
\end{przyklad}

\begin{przyklad}[Koequalizator w grupach]
Niech $f,g:A\to B$ będą homomorfizmami grup. Oznaczmy przez $N$ najmniejszą podgrupę normalną $B$ zawierającą elementy $f(a)g(a)^{-1}$ dla wszystkich $a\in A$. Wówczas projekcja $q:B\to B/N$ jest koequalizatorem $f,g$. Rzeczywiście $qf=qg$. Jeśli $v:B\to K$ spełnia $vf=vg$, to $f(a)g(a)^{-1}\in\Ker v$, a normalność jądra daje $N\subseteq\Ker v$; powstaje jedyny homomorfizm $\bar v:B/N\to K$ z $\bar vq=v$. W kategorii grup nie wystarcza więc zwykła relacja równoważności zbiorów: utożsamienia muszą być zgodne z mnożeniem.
\end{przyklad}


\subsection{Pullback: para spełniająca warunek zgodności}

\begin{definicja}[Kwadrat kartezjański, pullback]\label{def:pullback}
Niech $f:A\to C$ i $g:B\to C$ będą morfizmami kategorii $\mathcal C$. \emph{Pullbackiem} (iloczynem włóknistym) tej pary jest obiekt $P$ z morfizmami $p_A:P\to A$, $p_B:P\to B$ spełniającymi $f\circ p_A=g\circ p_B$, takimi że dla każdego $X$ i pary $u:X\to A$, $v:X\to B$ z $f\circ u=g\circ v$ istnieje dokładnie jeden $h:X\to P$ spełniający $p_A\circ h=u$, $p_B\circ h=v$. Diagram
\begin{center}
\begin{tikzpicture}[>=stealth]
 \node (a) at (0,2) {$P$};
 \node (b) at (4,2) {$B$};
 \node (c) at (0,0) {$A$};
 \node (d) at (4,0) {$C$};
 \draw[->] (a) -- node[above] {$p_B$} (b);
 \draw[->] (a) -- node[left] {$p_A$} (c);
 \draw[->] (b) -- node[right] {$g$} (d);
 \draw[->] (c) -- node[below] {$f$} (d);
\end{tikzpicture}
\end{center}
nazywamy wówczas \emph{kwadratem kartezjańskim}.
\end{definicja}

\begin{przyklad}[Iloczyn włóknisty zbiorów]\label{przyklad:włokna}
W $\Set$ konstrukcja ma bardzo konkretną postać:
\[
A\times_C B:=\{(a,b)\in A\times B:f(a)=g(b)\}.
\]
Rzutami są $p_A(a,b)=a$, $p_B(a,b)=b$. Jeśli $u:X\to A$ i $v:X\to B$ spełniają $fu=gv$, to jedyna możliwa funkcja $h$ jest dana przez $h(x)=(u(x),v(x))$. Warunek zgodności gwarantuje, że wartości $h$ rzeczywiście należą do $A\times_C B$.
\end{przyklad}

\begin{przyklad}[Włókno funkcji]
Dla $f:A\to C$ i ustalonego $c\in C$ weźmy $B=\{*\}$ i $g(*)=c$. Wówczas $A\times_C\{*\}$ można utożsamić ze zbiorem $f^{-1}(\{c\})$. A zatem \emph{włókno nad punktem} jest szczególnym przypadkiem pullbacku. Ten sam pomysł działa dla wiązek w topologii i geometrii różniczkowej.
\end{przyklad}

\begin{przyklad}[Ograniczenie rodziny do podzbioru]
Niech $p:E\to B$ będzie ciągłą funkcją, a $i:U\hookrightarrow B$ włączeniem podprzestrzeni. Pullback w $\Top$ jest $E\times_B U\cong p^{-1}(U)$ z topologią podprzestrzeni, wyposażony w odwzorowanie do $U$. Istotnie, rzut $(e,u)\mapsto e$ ma odwrotność $e\mapsto(e,p(e))$ na $p^{-1}(U)$; obie funkcje są ciągłe, bo ograniczenie $p:p^{-1}(U)\to U$ jest ciągłe dla topologii podprzestrzeni. Własność uniwersalna wynika z ciągłości parowania i ograniczenia przeciwdziedziny do podprzestrzeni. W języku wiązek oznacza to \emph{ograniczenie obiektu nad $B$ do $U$}.
\end{przyklad}

\begin{przyklad}[Wspólny fragment dwóch odwzorowań]
Niech $A=B=\mathbb R$, $C=\mathbb R$ oraz $f(a)=a^2$, $g(b)=b^2$. Pullback w $\Set$ jest zbiorem $\{(a,b):a^2=b^2\}$, czyli sumą prostych $b=a$ oraz $b=-a$ w $\mathbb R^2$. Nie jest nim zwykły iloczyn $\mathbb R^2$: warunek wspólnego obrazu usuwa większość par.
\end{przyklad}


\begin{przyklad}[Pullback jako konkretne równanie liniowe]
Pracujmy w $\Vect_{\mathbb R}$. Niech $f:\mathbb R^2\to\mathbb R$ będzie dane przez $f(x,y)=x+y$, a $g:\mathbb R\to\mathbb R$ przez $g(t)=2t$. Iloczyn włóknisty to podprzestrzeń
\[
P=\{(x,y,t)\in\mathbb R^3:x+y=2t\}=\{(x,2t-x,t):x,t\in\mathbb R\}.
\]
W szczególności $\dim P=2$. Jeżeli $u:Z\to\mathbb R^2$ i $v:Z\to\mathbb R$ są liniowe i spełniają $fu=gv$, to odwzorowanie $z\mapsto(u(z),v(z))$ jest liniowe, przyjmuje wartości w $P$ i jest jedynym odwzorowaniem zgodnym z rzutami. Obliczyliśmy więc zarówno przestrzeń, jak i jej własność uniwersalną.
\end{przyklad}

\begin{przyklad}[Przeciwobrazy jako pullbacki]
Dla funkcji $f:X\to Y$ oraz włączenia $i:B\hookrightarrow Y$ produkt włóknisty $X\times_Y B$ jest naturalnie izomorficzny ze zbiorem $f^{-1}(B)$. Ustalmy $f:\mathbb Z\to\mathbb Z$, $f(n)=n^2$, i $B=\{0,1,4\}$. Wówczas $f^{-1}(B)=\{-2,-1,0,1,2\}$, a rzut na $X$ identyfikuje produkt włóknisty z tym podzbiorem. Ten sposób rozumowania działa również dla podprzestrzeni topologicznych oraz dla przeciwobrazów podrozmaitości, gdy odpowiednie struktury istnieją.
\end{przyklad}


\subsection{Pushout: sklejanie po wspólnej części}

\begin{definicja}[Pushout]\label{def:pushout}
Niech $f:C\to A$ i $g:C\to B$ będą morfizmami kategorii $\mathcal C$. \emph{Pushoutem} pary $(f,g)$ nazywamy obiekt $Q$ z morfizmami $i_A:A\to Q$, $i_B:B\to Q$ spełniającymi $i_A\circ f=i_B\circ g$, takimi że dla każdego $X$ oraz $u:A\to X$, $v:B\to X$ z $u\circ f=v\circ g$ istnieje dokładnie jeden $h:Q\to X$ z $h\circ i_A=u$, $h\circ i_B=v$. Odpowiedni kwadrat
\begin{center}
\begin{tikzpicture}[>=stealth]
 \node (a) at (0,2) {$C$};
 \node (b) at (4,2) {$B$};
 \node (c) at (0,0) {$A$};
 \node (d) at (4,0) {$Q$};
 \draw[->] (a) -- node[above] {$g$} (b);
 \draw[->] (a) -- node[left] {$f$} (c);
 \draw[->] (b) -- node[right] {$i_B$} (d);
 \draw[->] (c) -- node[below] {$i_A$} (d);
\end{tikzpicture}
\end{center}
nazywamy \emph{kokartezjańskim}.
\end{definicja}

\begin{przyklad}[Sklejanie zbiorów]
W $\Set$ pushout $A\sqcup_C B$ ma postać $(A\sqcup B)/{\sim}$, gdzie $\sim$ jest najmniejszą relacją równoważności identyfikującą $j_A(f(c))$ z $j_B(g(c))$ dla każdego $c\in C$, a $j_A,j_B$ są zanurzeniami do sumy rozłącznej. Jeżeli $q:A\sqcup B\to Q$ jest rzutem ilorazowym, morfizmami pushoutu są $i_A=q\circ j_A$ i $i_B=q\circ j_B$. Dla zgodnej pary $u,v$ morfizm $[u,v]$ z sumy rozłącznej identyfikuje wszystkie wskazane pary, więc faktoryzuje się jednoznacznie przez $q$ według własności koequalizatora. Otrzymujemy wymagany morfizm $h:Q\to X$.
\end{przyklad}

\begin{przyklad}[Sklejanie przedziałów w topologii]
Weźmy dwie kopie przedziału $[0,1]$ i jednopunktowy zbiór $C=\{*\}$, który odwzorowujemy w punkt $1$ pierwszego przedziału i $0$ drugiego. Pushout w $\Top$ jest przestrzenią ilorazową otrzymaną przez sklejenie końców; jest homeomorficzny z $[0,2]$. Szczegółowy dowód homeomorfizmu podajemy poniżej. Wersja z identyfikacją obu końców \emph{jednego} przedziału prowadzi z kolei do okręgu $S^1$ i może być opisana przez odpowiedni koequalizator.
\end{przyklad}

\begin{uwaga}
Pullback to granica diagramu $A\to C\leftarrow B$, a pushout to kogranica diagramu $A\leftarrow C\to B$. Związek ten uzasadnimy ogólnie w sekcji~\ref{sec:granice}. Nazwy angielskie są często używane również w polskich tekstach, więc warto je rozpoznawać.
\end{uwaga}


\begin{przyklad}[Pushout dwóch homomorfizmów liczb całkowitych]
W kategorii $\Ab$ weźmy $C=A=B=\mathbb Z$ i morfizmy $f(n)=2n$, $g(n)=3n$. Pushout jest ilorazem
\[
Q=(\mathbb Z\oplus\mathbb Z)/\langle(2,-3)\rangle.
\]
Mapy kanoniczne są określone przez $i_A(a)=[(a,0)]$ oraz $i_B(b)=[(0,b)]$. Ponieważ $2i_A(1)=3i_B(1)$, diagram komutuje. Każda para homomorfizmów $u,v:\mathbb Z\to H$ do grupy abelowej $H$ z $u(2n)=v(3n)$ definiuje homomorfizm $\widehat h:\mathbb Z\oplus\mathbb Z\to H$ dany przez $(a,b)\mapsto u(a)+v(b)$; ma on w jądrze $(2,-3)$, więc jednoznacznie przechodzi na $Q$. Ponadto $Q\cong\mathbb Z$: funkcja $(a,b)\mapsto3a+2b$ jest surjektywna, ponieważ przyjmuje wartość $1$ na $(1,-1)$. Jeśli $3a+2b=0$, to $a$ jest parzyste, więc $a=2k$ i $b=-3k$. Odwrotnie każda para $(2k,-3k)$ należy do jądra. Pierwsze twierdzenie o izomorfizmie daje żądany izomorfizm. Zwykłe sklejanie zbiorów nie pozwoliłoby uzyskać tej własności addytywnej.
\end{przyklad}

\begin{przyklad}[Sklejenie dwóch odcinków i jego topologia]
Dla dwóch kopii $[0,1]$ ze sklejanymi końcami $(1,\mathrm{lewy})\sim(0,\mathrm{prawy})$ otrzymujemy przestrzeń $Q=([0,1]\sqcup[0,1])/{\sim}$ z topologią ilorazową. Funkcja określona na sumie rozłącznej wzorami $t\mapsto t$ na pierwszej kopii i $t\mapsto1+t$ na drugiej jest ciągła i stała na klasach relacji. Indukuje bijekcję $Q\to[0,2]$. Jest to homeomorfizm, gdyż $Q$ jest zwarty, a $[0,2]$ jest Hausdorffa. Ta konstrukcja realizuje pushout w $\Top$, a nie tylko w $\Set$, ponieważ topologia ilorazowa zapewnia ciągłość każdego odwzorowania określonego zgodnie na obu odcinkach.
\end{przyklad}

\subsection{Konstrukcje z iloczynów i equalizatorów}

\begin{stwierdzenie}\label{cat:prop:pullback-equalizer}
Jeżeli istnieje iloczyn $A\times B$ oraz equalizator $e:P\to A\times B$ pary $f\circ p_A,g\circ p_B:A\times B\to C$, to $P$ z rzutami $p_A\circ e,p_B\circ e$ jest iloczynem włóknistym $f,g$. Dualnie pushout $f:C\to A$, $g:C\to B$ można skonstruować jako koequalizator pary $j_A\circ f,j_B\circ g:C\to A\sqcup B$.
\end{stwierdzenie}
\begin{proof}
Równość po złożeniu z $e$ daje przemienność kwadratu. Dla zgodnej pary $u:X\to A$, $v:X\to B$ jedyny morfizm $\langle u,v\rangle:X\to A\times B$ spełnia $(f\circ p_A)\circ\langle u,v\rangle=(g\circ p_B)\circ\langle u,v\rangle$. Faktoryzacja przez $e$ daje morfizm do $P$ o wymaganych rzutach. Każdy inny taki morfizm po złożeniu z $e$ ma te same rzuty, więc daje $\langle u,v\rangle$ i jest równy pierwszemu przez jednoznaczność equalizatora. Dowód dla pushoutu otrzymujemy w kategorii przeciwnej, w której koprodukt jest iloczynem, a koequalizator equalizatorem.
\end{proof}

\begin{przyklad}[Moduły]\label{cat:ex:module-pull-push}
Dla lewych $R$-modułów i homomorfizmów $f:A\to C$, $g:B\to C$ pullback jest jądrem homomorfizmu $A\oplus B\to C$, $(a,b)\mapsto f(a)-g(b)$, z ograniczeniami rzutów. Dla $f:C\to A$, $g:C\to B$ pushout jest modułem
\[
(A\oplus B)/\operatorname{im}(f,-g)
\]
z morfizmami $a\mapsto[(a,0)]$, $b\mapsto[(0,b)]$. W obu przypadkach stwierdzenie~\ref{cat:prop:pullback-equalizer} stosujemy do equalizatora lub koequalizatora, które są odpowiednio jądrem różnicy i kojądrem różnicy. Istnienie tych modułów i własność ilorazu pochodzą z \atlasref{AA-16::aa:tw:modules-first-isomorphism}{AA-16: jądro, obraz i iloraz modułów}. Dla pushoutu relacja $[(f(c),0)]=[(0,g(c))]$ zachodzi, gdyż różnica tych par należy do podmodułu, przez który dzielimy.
\end{przyklad}

\begin{przyklad}[Topologia]\label{cat:ex:top-pull-push}
W $\Top$ equalizator otrzymujemy jako podprzestrzeń zbioru zgodności $E=\{a:f(a)=g(a)\}$; odwzorowanie o obrazie w $E$ jest ciągłe do $E$ wtedy i tylko wtedy, gdy jest ciągłe do $A$. Koequalizator jest ilorazem zbioru z topologią ilorazową, dla której $\bar v$ jest ciągłe wtedy i tylko wtedy, gdy $\bar v\circ q$ jest ciągłe (\atlasref{GT-04::thm:topog-iloraz-uniwersalnosc}{GT-04: własność uniwersalna ilorazu}). Stąd, na mocy stwierdzenia~\ref{cat:prop:pullback-equalizer}, pullback ma topologię podprzestrzeni iloczynu, a pushout topologię ilorazową sumy rozłącznej. Nie zakładamy, że przestrzenie ilorazowe są Hausdorffa.
\end{przyklad}

\subsection{Sklejanie kwadratów i stabilność morfizmów}

\begin{twierdzenie}[Lemat o sklejaniu kwadratów kartezjańskich]\label{cat:thm:pullback-pasting}
Niech poniższy diagram będzie przemienny i niech prawy kwadrat będzie kartezjański.
\begin{center}
\begin{tikzpicture}[>=stealth]
 \node (p) at (0,2) {$P$};
 \node (q) at (3.4,2) {$Q$};
 \node (r) at (6.8,2) {$R$};
 \node (x) at (0,0) {$X$};
 \node (y) at (3.4,0) {$Y$};
 \node (z) at (6.8,0) {$Z$};
 \draw[->] (p) -- node[above] {$a$} (q);
 \draw[->] (q) -- node[above] {$b$} (r);
 \draw[->] (p) -- node[left] {$p$} (x);
 \draw[->] (q) -- node[right] {$q$} (y);
 \draw[->] (r) -- node[right] {$r$} (z);
 \draw[->] (x) -- node[below] {$u$} (y);
 \draw[->] (y) -- node[below] {$v$} (z);
\end{tikzpicture}
\end{center}
Wówczas lewy kwadrat jest kartezjański wtedy i tylko wtedy, gdy zewnętrzny prostokąt jest kartezjański.
\end{twierdzenie}
\begin{proof}
Załóżmy najpierw kartezjańskość lewego kwadratu. Dla $s:W\to X$, $t:W\to R$ spełniających $v\circ u\circ s=r\circ t$, prawy kwadrat daje jedyny $k:W\to Q$ z $q\circ k=u\circ s$ i $b\circ k=t$. Lewy kwadrat daje jedyny $h:W\to P$ z $p\circ h=s$ i $a\circ h=k$. Wtedy $b\circ a\circ h=t$. Jeżeli $h'$ spełnia warunki zewnętrznego prostokąta, to $a\circ h'$ spełnia warunki na $k$, zatem $a\circ h'=k$; jedyność w lewym kwadracie daje $h'=h$.

Odwrotnie załóżmy kartezjańskość prostokąta. Dla $s:W\to X$, $k:W\to Q$ z $u\circ s=q\circ k$ para $s,b\circ k$ jest zgodna względem $v\circ u,r$. Prostokąt daje jedyny $h:W\to P$ z $p\circ h=s$ i $b\circ a\circ h=b\circ k$. Mamy $q\circ a\circ h=u\circ p\circ h=u\circ s=q\circ k$. Jednoznaczność w prawym kwadracie daje $a\circ h=k$. Każde inne rozwiązanie dla lewego kwadratu spełnia również warunki prostokąta, więc jest równe $h$.
\end{proof}

\begin{uwaga}
W kategorii przeciwnej i po odwróceniu kolejności kolumn otrzymujemy wersję dualną: jeżeli lewy kwadrat jest kokartezjański, to prawy jest kokartezjański wtedy i tylko wtedy, gdy kokartezjański jest zewnętrzny prostokąt.
\end{uwaga}

\begin{stwierdzenie}[Stabilność monomorfizmów i epimorfizmów]\label{cat:prop:base-change-mono}
W kwadracie kartezjańskim z definicji~\ref{def:pullback}, jeżeli $g:B\to C$ jest monomorfizmem, to $p_A:P\to A$ jest monomorfizmem. Dualnie w kwadracie kokartezjańskim, jeżeli $g:C\to B$ jest epimorfizmem, to $i_A:A\to Q$ jest epimorfizmem.
\end{stwierdzenie}
\begin{proof}
Niech $s,t:W\to P$ spełniają $p_A\circ s=p_A\circ t$. Z przemienności wynika
$g\circ p_B\circ s=f\circ p_A\circ s=f\circ p_A\circ t=g\circ p_B\circ t$.
Moniczność $g$ daje $p_B\circ s=p_B\circ t$. Jednoznaczność w iloczynie włóknistym daje $s=t$. W wersji dualnej dla $s,t:Q\to W$ z $s\circ i_A=t\circ i_A$ mamy $s\circ i_B\circ g=t\circ i_B\circ g$. Epiczność $g$ daje $s\circ i_B=t\circ i_B$, a własność pushoutu daje $s=t$.
\end{proof}

\begin{stwierdzenie}[Cofnięcie izomorfizmu]\label{cat:prop:base-change-iso}
Jeżeli w tym samym kwadracie kartezjańskim $g$ jest izomorfizmem, to $p_A$ jest izomorfizmem.
\end{stwierdzenie}
\begin{proof}
Para $\id_A,g^{-1}\circ f$ jest zgodna, więc daje morfizm $s:A\to P$ z $p_A\circ s=\id_A$ oraz $p_B\circ s=g^{-1}\circ f$. Morfizm $p_A$ jest monomorfizmem na mocy stwierdzenia~\ref{cat:prop:base-change-mono}. Z $p_A\circ s\circ p_A=p_A$ wynika $s\circ p_A=\id_P$, więc $s=p_A^{-1}$. Dualnie wypchnięcie izomorfizmu w kwadracie kokartezjańskim jest izomorfizmem.
\end{proof}

\subsection{Iloczyny włókniste rozmaitości}

\begin{przyklad}[Transwersalność i istnienie iloczynu włóknistego]\label{cat:ex:manifold-pullback}
Dla gładkich rozmaitości bez brzegu $M,N,S$ oraz gładkich odwzorowań $f:M\to S$, $g:N\to S$ załóżmy, że dla każdej pary $f(x)=g(y)=z$ zachodzi
\[
\mathrm df_x(T_xM)+\mathrm dg_y(T_yN)=T_zS.
\]
\atlasref{DG-15::thm:transwersalne-przeciecie}{DG-15: twierdzenie o transwersalnym iloczynie włóknistym} daje na $P=\{(x,y):f(x)=g(y)\}$ strukturę podrozmaitości zanurzonej w $M\times N$, o wymiarze $\dim M+\dim N-\dim S$ (jeśli $P$ jest niepuste). Jest to także pullback w kategorii gładkich rozmaitości: zgodna para gładkich odwzorowań daje gładkie odwzorowanie do $M\times N$ o obrazie zawartym w $P$. Jest ono gładkie jako odwzorowanie do $P$, ponieważ w lokalnych mapach dostosowanych do podrozmaitości składowe normalne są zerowe, a pozostałe składowe gładkie. Jednoznaczność wynika z rzutów.

Bez transwersalności pullback może nie istnieć. Dla $f=g:\mathbb R\to\mathbb R$, $f(t)=t^2$, przypuśćmy, że istnieje pullback $P$ z rzutami $p,q:P\to\mathbb R$. Z własności uniwersalnej dla punktu istnieje jedyny $o\in P$ o $p(o)=q(o)=0$. Zgodne pary $(t,t)$ i $(t,-t)$ dają gładkie krzywe $\gamma_+,\gamma_-:\mathbb R\to P$ przechodzące przez $o$. Dla $F=(p,q):P\to\mathbb R^2$ mamy
\[
\mathrm dF_o(\gamma'_+(0))=(1,1),\qquad
\mathrm dF_o(\gamma'_-(0))=(1,-1).
\]
Te wektory generują $\mathbb R^2$, więc $\mathrm dF_o$ jest surjektywna. Z twierdzenia o submersji obraz pewnego otoczenia $o$ zawiera otoczenie $(0,0)$ w $\mathbb R^2$. Jednak z przemienności $p^2=q^2$ obraz $F$ zawiera się w sumie dwóch prostych $y=x$ i $y=-x$, która nie zawiera żadnego otoczenia $(0,0)$. Sprzeczność. Sam zbiór zgodności nie zapewnia więc istnienia iloczynu włóknistego w kategorii rozmaitości.
\end{przyklad}

